\subsubsection*{Численные примеры}

\begin{idea}[title={Откуда этот материал}]
На лекции симплекс-метод и метод искусственного базиса разбирались только в
общем виде. Ниже оба алгоритма прогнаны на числах — по тем же формулам, что в
\S\ref{subsec:simplex-scheme} и \S\ref{subsec:artificial-basis}.
\end{idea}

\begin{example}[Симплекс-метод]\label{ex:simplex-run}
\[
  \begin{cases}
    F(x) = -3x_1 - x_2 \to \min , \\
    x_1 + x_2 + x_3 = 4 , \\
    2x_1 + x_2 + x_4 = 5 , \\
    x_j \ge 0 , \quad j = \overline{1,4} .
  \end{cases}
\]
Здесь $m = 2$, $n = 4$,
\[
  A[M,N] = \begin{pmatrix} 1 & 1 & 1 & 0 \\ 2 & 1 & 0 & 1 \end{pmatrix} ,
  \qquad
  b[M] = \begin{pmatrix} 4 \\ 5 \end{pmatrix} ,
  \qquad
  c\T[N] = (-3, -1, 0, 0) .
\]
Переменные $x_3, x_4$ вошли в задачу как балансовые, поэтому начальный базис
виден сразу: $N_0 = \{3, 4\}$, $x^0[N] = (0, 0, 4, 5)\T$, $F(x^0) = 0$.

\medskip
\noindent\textbf{Итерация 0.} Так как $A[M,N_0] = E$, то и $B[N_0,M] = E$.
Потенциалы находятся из системы
\[
  A\T[M,N_0]\,y^0[M] = c[N_0] = (0,0)\T
  \qquad\Longrightarrow\qquad
  y^0[M] = (0,0)\T .
\]
Оценки небазисных индексов $L_0 = \{1,2\}$:
\[
  d[1] = c[1] - (y^0)\T a_1 = -3 - 0 = -3 , \qquad
  d[2] = c[2] - (y^0)\T a_2 = -1 - 0 = -1 .
\]
Обе отрицательны; берём наиболее отрицательную — $j_0 = 1$. Направление:
\[
  u^0[N_0] = B[N_0,M]\,a_1[M] = \begin{pmatrix} 1 \\ 2 \end{pmatrix} ,
  \qquad u^0[1] = -1 , \qquad u^0[2] = 0 ,
\]
то есть $u^0[3] = 1$, $u^0[4] = 2$. Шаг по правилу минимального отношения:
\[
  \theta_0 = \min\Big\{ \frac{x^0[3]}{u^0[3]},\ \frac{x^0[4]}{u^0[4]} \Big\}
           = \min\Big\{ \frac{4}{1},\ \frac{5}{2} \Big\} = \frac52 ,
  \qquad i_0 = 4 .
\]
Новый опорный вектор $x^1[N] = x^0[N] - \theta_0 u^0[N]$:
\[
  x^1[1] = 0 - \tfrac52(-1) = \tfrac52 , \quad
  x^1[2] = 0 , \quad
  x^1[3] = 4 - \tfrac52 = \tfrac32 , \quad
  x^1[4] = 5 - \tfrac52\cdot 2 = 0 ,
\]
\[
  x^1[N] = \big(\tfrac52,\ 0,\ \tfrac32,\ 0\big)\T ,
  \qquad N_1 = \{3, 1\} ,
  \qquad F(x^1) = -\tfrac{15}{2} .
\]
Контроль: $F(x^1) = F(x^0) + \theta_0 d[j_0] = 0 + \tfrac52 \cdot (-3)
= -\tfrac{15}{2}$ — совпадает.

\medskip
\noindent\textbf{Итерация 1.} Базисные столбцы $a_3 = (1,0)\T$, $a_1 = (1,2)\T$:
\[
  A[M,N_1] = \begin{pmatrix} 1 & 1 \\ 0 & 2 \end{pmatrix} ,
  \qquad
  B[N_1,M] = \begin{pmatrix} 1 & -\tfrac12 \\[2pt] 0 & \tfrac12 \end{pmatrix} .
\]
Потенциалы из $A\T[M,N_1]\,y^1[M] = c[N_1] = (0, -3)\T$:
\[
  \begin{cases}
    y^{(1)} = 0 , \\
    y^{(1)} + 2 y^{(2)} = -3 ,
  \end{cases}
  \qquad\Longrightarrow\qquad
  y^1[M] = \big(0,\ -\tfrac32\big)\T .
\]
Оценки по $L_1 = \{2,4\}$:
\[
  d[2] = -1 - \big(0\cdot 1 + (-\tfrac32)\cdot 1\big) = \tfrac12 \ge 0 ,
  \qquad
  d[4] = 0 - \big(0\cdot 0 + (-\tfrac32)\cdot 1\big) = \tfrac32 \ge 0 .
\]
Все оценки неотрицательны, значит по теореме~\ref{thm:simplex-opt}
\[
  x_*[N] = \big(\tfrac52,\ 0,\ \tfrac32,\ 0\big)\T ,
  \qquad F(x_*) = -\tfrac{15}{2} .
\]
\end{example}

\begin{idea}
Проверка геометрией: в исходных переменных задача — это
$3x_1 + x_2 \to \max$ при $x_1 + x_2 \le 4$, $2x_1 + x_2 \le 5$, $x \ge 0$.
Вершины допустимого многоугольника: $(0,0)$, $(0,4)$, $(1,3)$, $(2{,}5;\,0)$ со
значениями $0$, $4$, $6$, $7{,}5$. Максимум $7{,}5$ в точке $(2{,}5;\,0)$ — то
же, что дал метод за один шаг.
\end{idea}

\begin{example}[Метод искусственного базиса]\label{ex:artificial-run}
Требуется найти допустимое решение системы
\[
  \begin{cases}
    x_1 + x_2 = 3 , \\
    x_1 - x_2 = 1 , \\
    x_1, x_2 \ge 0 .
  \end{cases}
\]
Очевидного базиса нет. Правая часть уже неотрицательна, поэтому сразу вводим
искусственные переменные $y_1, y_2$ (индексы $3$ и $4$) и решаем вспомогательную
задачу
\[
  \begin{cases}
    W = y_1 + y_2 \to \min , \\
    x_1 + x_2 + y_1 = 3 , \\
    x_1 - x_2 + y_2 = 1 , \\
    x_1, x_2, y_1, y_2 \ge 0 ,
  \end{cases}
  \qquad
  A'[M,N] = \begin{pmatrix} 1 & 1 & 1 & 0 \\ 1 & -1 & 0 & 1 \end{pmatrix} ,
  \quad {c'}\T[N] = (0,0,1,1) .
\]
Старт: $N_0 = \{3,4\}$, $x^0[N] = (0,0,3,1)\T$, $W = 4$.

\medskip
\noindent\textbf{Итерация 0.} $A'[M,N_0] = E$, $B = E$, $y^0[M] = c'[N_0] = (1,1)\T$.
\[
  d[1] = 0 - (1\cdot 1 + 1\cdot 1) = -2 , \qquad
  d[2] = 0 - \big(1\cdot 1 + 1\cdot(-1)\big) = 0 .
\]
Выбираем $j_0 = 1$. Направление $u^0[N_0] = B\,a_1 = (1,1)\T$, то есть
$u^0[3] = u^0[4] = 1$, $u^0[1] = -1$, $u^0[2] = 0$. Шаг
\[
  \theta_0 = \min\Big\{ \frac31,\ \frac11 \Big\} = 1 , \qquad i_0 = 4
  \quad (\text{из базиса выходит } y_2) .
\]
\[
  x^1[N] = (1,\ 0,\ 2,\ 0)\T , \qquad N_1 = \{3,1\} , \qquad W = 2 .
\]

\medskip
\noindent\textbf{Итерация 1.} Базисные столбцы $a_3 = (1,0)\T$, $a_1 = (1,1)\T$:
\[
  A'[M,N_1] = \begin{pmatrix} 1 & 1 \\ 0 & 1 \end{pmatrix} ,
  \qquad
  B[N_1,M] = \begin{pmatrix} 1 & -1 \\ 0 & 1 \end{pmatrix} ,
  \qquad
  c'[N_1] = (1, 0)\T .
\]
Потенциалы: $y^{(1)} = 1$, $y^{(1)} + y^{(2)} = 0 \implies y^1[M] = (1,-1)\T$.
Оценки по $L_1 = \{2,4\}$:
\[
  d[2] = 0 - \big(1\cdot 1 + (-1)\cdot(-1)\big) = -2 < 0 ,
  \qquad
  d[4] = 1 - \big(1\cdot 0 + (-1)\cdot 1\big) = 2 .
\]
Выбираем $j_1 = 2$:
\[
  u^1[N_1] = B[N_1,M]\,a_2[M]
    = \begin{pmatrix} 1 & -1 \\ 0 & 1 \end{pmatrix}
      \begin{pmatrix} 1 \\ -1 \end{pmatrix}
    = \begin{pmatrix} 2 \\ -1 \end{pmatrix} ,
\]
то есть $u^1[3] = 2$, $u^1[1] = -1$, $u^1[2] = -1$, $u^1[4] = 0$. Положительна
только компонента при $i = 3$, поэтому
\[
  \theta_1 = \frac{x^1[3]}{u^1[3]} = \frac22 = 1 , \qquad i_1 = 3
  \quad (\text{из базиса выходит } y_1) .
\]
\[
  x^2[N] = (2,\ 1,\ 0,\ 0)\T , \qquad N_2 = \{1,2\} , \qquad W = 0 .
\]
Минимум вспомогательной задачи равен нулю, искусственные переменные вышли из
базиса, значит по теореме~\ref{thm:artificial-basis} вектор $(2,1)\T$ допустим в
исходной системе: $2 + 1 = 3$, $2 - 1 = 1$. Базис $N_2 = \{1,2\}$ и берётся
начальным для решения самой задачи обычным симплекс-методом.
\end{example}

\begin{idea}
Обе искусственные переменные вышли из базиса сами, поэтому шаг~8 алгоритма
(замена столбцов) не понадобился. Так бывает не всегда: если бы система
содержала зависимую строку, одна из $y_i$ осталась бы в базисе с нулевым
значением.
\end{idea}
